\section{Related Work}
\noindent\textbf{Speech-Driven 3D Facial Animation.}
Speech-driven 3D facial animation seeks to generate temporally synchronized facial motions from audio. Early approaches formulated the task as direct regression from acoustic features to facial geometry. VOCA~\cite{cudeiro2019voca} predicts speaker-specific 3D vertices, while MeshTalk~\cite{richard2021meshtalk} introduces a categorical latent space for non-verbal motions. Transformer-based models later improved long-range temporal modeling: FaceFormer~\cite{fan2022faceformer} employs an autoregressive transformer, and UniTalker~\cite{fan2024unitalker} enables unified training across datasets.
A key challenge is the inherent one-to-many mapping from audio to facial motion, which often leads to over-smoothed results. To address this, CodeTalker~\cite{xing2023codetalker} incorporates discrete motion priors via vector quantization. Diffusion-based methods such as SAiD~\cite{park2023said} and FaceDiffuser~\cite{stan2023facediffuser} better capture motion diversity and stochasticity. In production settings, NVIDIA Audio2Face-3D~\cite{chung2025audio2face} provides a robust real-time pipeline, while SentiAvatar~\cite{jin2026sentiavatar} introduces planning strategies for improved semantic consistency at the sentence level. More recently, Ex-Omni~\cite{zhang2026ex}, where the omni-modal mainly provides semantic understanding and speech-unit representations, employs a separate lightweight facial generator that maps these units to blendshape coefficients.
Building on these advances, AudioMouth formulates mouth-motion prediction as fixed-schema coefficient generation with a multimodal language model conditioned on audio and explicit linguistic inputs. We complement token-level supervision with sequence-level refinement that evaluates the generated trajectories in coefficient, temporal, and geometric terms.

\noindent\textbf{Facial Motion Representations.}
The choice of facial motion representation significantly impacts controllability, interpretability, and compatibility with modern multimodal models. Vertex-based representations~\cite{cudeiro2019voca,fan2022faceformer} preserve fine geometric details but lack semantic interpretability. Parametric models such as FLAME~\cite{li2017learning} achieve compactness via PCA, yet their latent dimensions are not naturally aligned with linguistic or phonetic concepts.
ARKit blendshapes~\cite{apple_arkit} provide semantically meaningful facial action controls, such as jaw opening, cheek puffing, and mouth closure, while SemanticFace~\cite{kang2026semanticface} demonstrates the benefit of language-aligned supervision for estimating these coefficients from images.
AudioMouth adopts 32 mouth-related ARKit channels as an interpretable output space that connects articulatory cues with explicit facial controls. This coefficient space also supports SLMD, which measures the resulting mouth deformation through a fixed blendshape basis.

\noindent\textbf{Linguistic and Articulatory Priors.}
Speech articulation is fundamentally tied to phonetic and linguistic structure: bilabial consonants involve lip closure, open vowels generally require greater jaw opening, and phoneme transitions shape temporal dynamics~\cite{cohen1993modeling,taylor2012dynamic,edwards2016jali}. Classic and viseme-based approaches~\cite{cohen1993modeling,edwards2016jali,taylor2017deep,zhou2018visemenet,bao2023learning} explicitly leverage these relations to guide speech animation.
Recent works incorporate higher-level cues. Emotion- and sentiment-aware methods~\cite{peng2023emotalk,jin2026sentiavatar} add control signals beyond raw audio, demonstrating the value of non-acoustic information. Recent audio-text or linguistically guided methods~\cite{fan2022joint,xu2024kmtalk} further suggest that textual or linguistic cues can reduce the ambiguity of audio-to-motion generation.
AudioMouth brings these cues into temporally aligned training units containing audio, transcripts, phonemes, and target coefficients. Within each unit, explicit phoneme-to-articulation guidance conditions coefficient generation alongside the acoustic input, while a distribution-balanced objective addresses the imbalance in coefficient values.
